\documentclass[11pt]{article}
\usepackage[utf8]{inputenc}
\usepackage{amsmath,amssymb,amsthm,mathtools,booktabs,array}
\usepackage[a4paper,margin=27mm]{geometry}
\usepackage[hidelinks]{hyperref}
\usepackage{microtype}
\newcommand{\workbenchurl}{https://github.com/KokunoYumeto/zeta-and-l-functions-workbench}
\newcommand{\recordurl}{https://github.com/KokunoYumeto/mathematics-workbenches}
\newcommand{\Q}{\mathbb Q}
\newcommand{\C}{\mathbb C}
\newcommand{\Z}{\mathbb Z}
\newcommand{\Tr}{\operatorname{Tr}}
\newcommand{\N}{\operatorname{N}}
\newcommand{\Rea}{\operatorname{Re}}
\newcommand{\im}{\operatorname{Im}}
\newcommand{\ord}{\operatorname{ord}}
\newtheorem{theorem}{Theorem}[section]
\newtheorem{proposition}[theorem]{Proposition}
\newtheorem{lemma}[theorem]{Lemma}
\newtheorem{corollary}[theorem]{Corollary}
\theoremstyle{definition}
\newtheorem{definition}[theorem]{Definition}
\theoremstyle{remark}
\newtheorem{remark}[theorem]{Remark}
\title{CM Hecke $L$-functions after Alp\"oge:\\finite phases and a prime-ideal window criterion}
\author{}
\date{24 November 2026}
\begin{document}
\maketitle
\begin{abstract}
Alp\"oge's 2026 theorem proves the Riemann hypothesis for every Dirichlet
$L$-function, and his proof contains an abstract analytic criterion based on
one-sided bounds on long intervals.  We verify that this criterion applies
to every finite positive combination of unitary Hecke characters of an
imaginary quadratic field having nonzero angular infinity type.  Over
$\Q(\sqrt{-3})$, the geometric order-three automorphism of a curve with
$j=0$ commutes with Frobenius at every good prime ideal; even over $\Q$,
sixth powers remove its finite phase.  We distinguish this calculation from
a global bound for prime sums, identify the cubic twist between the curves
$y^2=x^3+1$ and $y^2=x^3+16$, and describe one precise prime-ideal inequality
that would prove the Riemann hypotheses for their Hecke characters and
their common cube.  We also continue the sextic probe's auxiliary Euler
product and identify its possible zero cancellation with a shifted Hecke
factor of angular type twelve.  The prime-ideal inequality is not proved
here.  The mathematics
and text were produced by GPT-6 Astra in ChatGPT Pro mode; the principal
analytic criterion was independently re-proved by another instance of the
same model, and no human has checked this text.
\end{abstract}

\section{Main results and their scope}

The first problem beyond Alp\"oge's Dirichlet theorem in this discussion is
the Riemann hypothesis for the CM Hecke characters of two elliptic curves,
not a new local bound for their Frobenius eigenvalues.  Alp\"oge's announcement
describes the sixfold metaplectic cover and the Eisenstein CM curve in its
third step and accompanying footnote \cite[p.~3, step~3 and footnote~3]{Alpoge}.
Its twenty-step outline ends with the Dirichlet theorem.  The published
formal conclusions quantify over Dirichlet characters of rational integer
moduli \cite[\texttt{DirichletLNonvanishingRightOfHalf}]{AlpogeCode}.
The Lean community's verification supports these conclusions
\cite{LeanCommunity,MITVerification}.  It does not change the domain of
their quantifiers.

The preceding cubic-theta argument of OpenAI gives a different useful
reference point.  Its Theorem~1.1, dated 30 September 2026, excludes
zeros with real part greater than $7/8$ for finite-order Hecke
characters of $\Q(\sqrt{-3})$ and for every Dirichlet character
\cite[p.~4]{OpenAIQRH}.  Hu's subsequent extension treats continuous
unitary Hecke characters, including nonzero infinity type
\cite[Theorem~1.2]{Hu}.  Section~6 below examines its principal Euler
factor and the later moment improvements.  These domains and
conclusions must be kept distinct when considering the CM targets.

Write $K=\Q(\omega)$, where $\omega^2+\omega+1=0$, and
$\mathcal O=\Z[\omega]$.  Let $\psi_{36}$ and $\psi_{27}$ denote the
algebraic Hecke characters of infinity type one of
\[
 E_{36}:y^2=x^3+1,\qquad E_{27}:y^2=x^3+16.
\]
The latter curve is isomorphic over $\Q$ to $y^2+y=x^3$: use
$X=4x$, $Y=8y+4$.  We use arithmetic normalization, so the conjectural
critical lines of $L(s,\psi_{36})$ and $L(s,\psi_{27})$ are
$\Rea s=1$.  Their common cube has critical line $\Rea s=2$.
For the analytic criterion we instead use the unitary characters
\begin{equation}\label{eq:targets}
 \lambda_1=\psi_{36}\N^{-1/2},\qquad
 \lambda_2=\psi_{27}\N^{-1/2},\qquad
 \lambda_3=\psi_{36}^{3}\N^{-3/2},
\end{equation}
whose common critical line is $\Rea s=1/2$.
All Hecke characters below are understood in their primitive form.

\begin{theorem}[Finite phases]\label{thm:phase}
Let $V$ be a finite-dimensional complex vector space, let
$F\in\operatorname{GL}(V)$, and let $\Gamma\subset\operatorname{GL}(V)$ be
a finite group normalized by $F$.  Let $h$ be the order of the automorphism
of $\Gamma$ induced by conjugation by $F$, and let $e$ be the exponent of
$\Gamma$.  For every $s\in\Gamma$,
\[
 (sF)^{he}=F^{he}.
\]
If $C,B>0$ and $|\Tr((sF)^{hen})|\leq C B^{hen}$ for every $n\geq1$, then every
eigenvalue of $F$ has modulus at most $B$.

On the cohomology of $y^2=x^3+b$, at a good prime ideal of $K$, the
automorphism $s(x,y)=(\omega x,y)$ commutes with Frobenius.
At a rational good prime $p\geq5$, the corresponding operators satisfy
$FsF^{-1}=s^p$ and $(sF)^6=F^6$.
These identities persist on every tensor power.
\end{theorem}

This theorem settles the finite-order phase question for the geometric CM
action.  It does not identify a prime sum with a finite-field trace, or
bound a sum over different rational primes.  The relevant global condition
is the following.

For an integral ideal $\mathfrak a$, set
$\Lambda_K(\mathfrak a)=\log\N\mathfrak p$ if
$\mathfrak a=\mathfrak p^r$ for a prime ideal and $r\geq1$, and set it to
zero otherwise.  Extend a primitive Hecke character by zero on ideals not
coprime to its conductor.  If $\lambda$ is unitary and has nonzero angular
infinity type, put
\begin{equation}\label{eq:Gdef}
 G_\lambda(x)=\Rea\left\{
 \sum_{\N\mathfrak a\leq x}\Lambda_K(\mathfrak a)\lambda(\mathfrak a)
       \left(\frac{x}{\N\mathfrak a}-1\right)
       +x\frac{L'(1,\lambda)}{L(1,\lambda)}\right\},\qquad x\geq1.
\end{equation}
The denominator in this definition is nonzero by Lemma~\ref{lem:one}.

\begin{theorem}[Hecke-character window criterion, after Alp\"oge]
\label{thm:window}
Let $K$ be an imaginary quadratic field.  Let $\lambda_1,\ldots,\lambda_r$
be primitive unitary Hecke characters of $K$, each with nonzero angular
infinity type, and let $a_1,\ldots,a_r$ be positive real numbers.  Define
$G(x)=\sum_j a_jG_{\lambda_j}(x)$.
Suppose that, for every $\theta\in(1/2,1)$, there exist
$\delta_\theta>0$ and $C_\theta\in\mathbb R$ such that, for every
$Y_0>1$, there exists $Y\geq Y_0$ with
\begin{equation}\label{eq:window}
 G(x)\leq C_\theta x^\theta
 \quad\hbox{for every }x\in[Y,Y^{1+\delta_\theta}].
\end{equation}
Then every nontrivial zero of every $L(s,\lambda_j)$ has real part $1/2$.
The proof invokes Alp\"oge's abstract localised Landau theorem
\cite[\texttt{Deduction.landau\_window}]{AlpogeWindow}; all of its
Hecke-character hypotheses are verified below.
\end{theorem}

\begin{corollary}\label{cor:three}
For the three characters in \eqref{eq:targets}, a single function
$G=G_{\lambda_1}+G_{\lambda_2}+G_{\lambda_3}$ satisfying
\eqref{eq:window} proves the Riemann hypotheses for both elliptic curves
and for their common cubic Hecke character.  Together with the CM
symmetric-cube decomposition, it proves the Riemann hypothesis for the
symmetric-cube $L$-function of $\eta(6\tau)^4$.
\end{corollary}

The corollary is a reduction, and the hypothesis \eqref{eq:window} remains
open.  No assertion of Hecke-character RH is made.  The content added here
to the analytic mechanism is the verification of its hypotheses for
positive infinity type, including positive sums of several characters.
The argument works for every fixed infinity type; it supplies no uniform
arithmetic bound as that type or its conductor grows.

\section{Proof overview}

There are two separate operations to justify.  First, a finite phase can
be removed from powers of Frobenius when Frobenius normalizes its finite
group.  This is a group calculation.  An elementary pole argument then
shows why genuine power traces, with their nonnegative integer
multiplicities, detect all eigenvalues.  In the CM example the necessary
normalization relation follows by raising the coordinates to their
residue-field cardinality.

Second, a global $L$-function is detected by a centered sum over all prime
ideals.  Taking a real part must be done holomorphically: one pairs the
character with its conjugate before taking a Mellin transform.  At a zero
the residues are positive weighted multiplicities divided by
$\rho(1-\rho)$.  Coincident zeros therefore reinforce one another; they
cannot disappear in the sum.  The centering term removes the pole at one.
These facts verify the abstract analytic theorem in Alp\"oge's proof.
What is still needed is a construction giving its one-sided prime-ideal
bound.  A finite-order action on an elliptic curve does not give that
bound by itself.

\section{Finite groups, Frobenius, and the cubic twist}

This section identifies exactly which phase relations the CM geometry
provides.

\begin{proof}[Proof of Theorem~\ref{thm:phase}]
Let $\varphi(\gamma)=F\gamma F^{-1}$.  Multiplying in the semidirect
product gives
\[
 (sF)^h=\bigl(s\varphi(s)\cdots\varphi^{h-1}(s)\bigr)F^h=tF^h,
 \qquad t\in\Gamma.
\]
Since $\varphi^h=1$, $F^h$ commutes with $t$.  Thus
$(sF)^{he}=t^eF^{he}=F^{he}$.

Put $N=he$ and let $\alpha_1,\ldots,\alpha_d$ be the eigenvalues of
$F$, with algebraic multiplicity.  The assumed trace bound says that
$\sum_{n\geq1}(\sum_i\alpha_i^{Nn})z^n$ is analytic for
$|z|<B^{-N}$.  Near zero this function equals
\[
 \sum_{i=1}^d\frac{\alpha_i^Nz}{1-\alpha_i^Nz}.
\]
For a fixed nonzero value $\beta=\alpha_i^N$, all coincident terms have
the same nonzero residue multiplied by a positive integer.  They cannot
cancel.  Hence this rational function has a pole at $\beta^{-1}$, which
cannot lie inside $|z|<B^{-N}$.  Thus $|\alpha_i|\leq B$.
No semisimplicity assumption on $F$ is required: traces of powers count
the eigenvalues with their algebraic multiplicities.

On the curve, the $q$-power Frobenius satisfies
$F s=s^q F$ by direct substitution in the two coordinates.  At a prime
ideal of $K$ not above three, the residue field contains $\mu_3$, so
$q\equiv1\pmod3$.  At a rational good prime, if $p\equiv1\pmod3$ then
$(sF)^3=F^3$.  If $p\equiv-1\pmod3$, then
$(sF)^2=s(FsF^{-1})F^2=ss^{-1}F^2=F^2$.
In both cases $(sF)^6=F^6$.  Functoriality gives the same identities on
cohomology and tensor powers.  Arithmetic and geometric Frobenius use
inverse conjugation exponents, which are identical here because the
exponents are $1$ and $-1$ modulo three.
\end{proof}

The relation $FsF^{-1}=s^{-1}$ in this proof is different from
$sFs^{-1}=F^{-1}$.  In the latter situation let
\[
 F=\begin{pmatrix}R&0\\0&R^{-1}\end{pmatrix},\qquad
 s=\begin{pmatrix}0&1\\1&0\end{pmatrix},\qquad R>1.
\]
Then $(sF)^2=1$, although $F$ has an eigenvalue $R$.
Here $F$ does not normalize the two-element group generated by $s$.
The normalizer hypothesis of Theorem~\ref{thm:phase} is therefore
substantive.

\begin{proposition}[Unit colour and infinity type]\label{prop:colour}
Suppose an algebraic Hecke character over $K$ has infinity type $m$ and
is written on principal ideals as
\[
 \psi((\alpha))=\varepsilon(\alpha)\alpha^m
 \quad ((\alpha,\mathfrak f)=1).
\]
Then $\varepsilon(u)=u^{-m}$ for every unit $u\in\mu_6$.
This colour determines $m$ modulo six, not its full infinity type.
In particular a cubic power of a type-one character is trivial on the
order-three part of the unit colour and nontrivial on its order-two part.
The centre of $\operatorname{GL}_2(\mathbb F_3)$ has no element of
order three.
\end{proposition}

\begin{proof}
The ideal $(u)$ is the unit ideal, so
$1=\varepsilon(u)u^m$.  This proves the first assertion and its
consequence for $m=3$.  Since $\mathcal O$ has class number one, the rule
$(\alpha)\mapsto\alpha^6$ is a conductor-one Hecke character: changing
a generator multiplies its value by $u^6=1$.  Twisting by this character
changes the infinity type from $m$ to $m+6$ and leaves the finite unit
colour unchanged.  It is not a finite-order twist.  For example, for
$\alpha=3+\omega$ of norm seven the normalized value is
$(\alpha/\bar\alpha)^3$; its valuations at the two primes above seven
are opposite and nonzero, so it cannot be a root of unity.
Finally, a matrix commuting with both elementary transvections in
$\operatorname{GL}_2(\mathbb F_3)$ is scalar.  The nonzero scalars are
$\pm1$, which proves the assertion about the centre.
\end{proof}

\begin{proposition}[The two CM curves and their common cube]
\label{prop:twist}
Let $\alpha^3=2$.  The map
\[
 E_{36}\longrightarrow E_{27},\qquad (x,y)\longmapsto(2\alpha x,4y)
\]
is an isomorphism over $K(\alpha)$.  Its descent cocycle takes values in
the cubic automorphisms.  Consequently the two CM characters differ by
a cubic Hecke character, and
$\psi_{36}^{3}=\psi_{27}^{3}$ as primitive characters, after matching
their chosen CM embeddings.

More explicitly, let $p\equiv1\pmod3$ be a good rational prime, let
$\chi$ be a cubic character of $\mathbb F_p^\times$, and let $\phi$ be
its quadratic character.  Extend both by zero at zero, and put
\[
 J=\sum_{t\in\mathbb F_p}\chi(t)\phi(1-t).
\]
Then, for $p\nmid b$,
\begin{equation}\label{eq:jacobi}
 a_p(y^2=x^3+b)=-\phi(b)
       \bigl(\chi(b)J+\chi(b)^2\bar J\bigr).
\end{equation}
In particular $b=16$ changes the Frobenius value $-J$ for $b=1$ by
the factor $\chi(2)$.
\end{proposition}

\begin{proof}
The coordinate map satisfies
$(4y)^2=(2\alpha x)^3+16$ since $\alpha^3=2$.
If $\sigma(\alpha)=\zeta\alpha$, its conjugate differs by
$(X,Y)\mapsto(\zeta X,Y)$, proving the cubic-twist assertion on the
one-dimensional CM summands of the Tate module.  Cubing removes this
cocycle.  The identity of characters then includes their primitive local
factors, not only the factors at good primes.

For completeness, count points directly.  The number of solutions of
$x^3=t$ is $1+\chi(t)+\chi(t)^2$, including $t=0$.  Therefore
\[
 -a_p=\sum_t(1+\chi(t)+\chi(t)^2)\phi(t+b).
\]
The term with coefficient one sums to zero.  Substitute $t=-bu$ in
the other two terms.  Since $\chi(-1)=1$, this gives
$-a_p=\phi(b)(\chi(b)J+\chi(b)^2\bar J)$.
For $b=16$, $\phi(16)=1$ and $\chi(16)=\chi(2)$.
At an inert good prime, $x\mapsto x^3$ is a permutation, so the same
point count gives $a_p=0$ for both curves.
\end{proof}

As a small example, at $p=7$ one choice of cubic character gives
$-J=-1+2\omega$, of trace $-4$ and norm seven.  Multiplication by
$\omega$ gives $-2-3\omega$, of trace $-1$, for the cubic twist.
Both cubes are equal, with trace twenty.  The exact program
\emph{CM phase checks}, available at \url{\workbenchurl}, verifies
the Jacobi identity and the cubic traces at every good prime below one
thousand.  These checks are finite corroboration of the preceding proofs.

\section{The analytic criterion for positive infinity type}

This section proves Theorem~\ref{thm:window} by checking the complete
analytic interface of Alp\"oge's localised Landau theorem.
We use the classical continuation and functional equation for Hecke
$L$-functions; a detailed adelic proof is in
\cite[Theorem~4 and its proof, pp.~10--11; Corollary~2, p.~9]{Prasad}.
For a unitary character with nonzero angular component, the finite
$L$-function is entire.  Its completed functional equation relates
$\lambda$ and $\bar\lambda$ across $\Rea s=1/2$.

\begin{lemma}\label{lem:one}
If $\lambda$ is unitary and has nonzero angular infinity type, then
$L(s,\lambda)\ne0$ on $\Rea s\geq1$.
\end{lemma}

\begin{proof}
Absolute convergence of the Euler product gives the assertion for
$\Rea s>1$.  To exclude a zero at $1+it$, replace $\lambda$ by
$\lambda\N^{-it}$, so that it suffices to work at one.  This replacement
and its square still have nonzero angular type, hence entire
$L$-functions.  For real $\sigma>1$ the Euler products imply
\[
 \zeta_K(\sigma)^3|L(\sigma,\lambda)|^4
                      |L(\sigma,\lambda^2)|\geq1.
\]
Indeed, at a prime unramified for $\lambda$, each term in the logarithm
is a positive multiple of
$3+4\cos t+\cos(2t)=2(1+\cos t)^2\geq0$.
At a prime ramified for $\lambda$, the contribution is either three or
$3+\cos t$, again nonnegative; $\lambda^2$ here denotes its primitive
character.  If $L(1,\lambda)$ had a zero of order at least one, the left
side would tend to zero as $\sigma\downarrow1$, since $\zeta_K$ has a
simple pole and $L(s,\lambda^2)$ is bounded at one.  This is a
contradiction.
\end{proof}

\begin{proof}[Proof of Theorem~\ref{thm:window}]
Define the meromorphic function
\[
 D(s)=\frac12\sum_{j=1}^r a_j
 \left(\frac{L'(s,\lambda_j)}{L(s,\lambda_j)}+
       \frac{L'(s,\bar\lambda_j)}{L(s,\bar\lambda_j)}\right).
\]
It satisfies $D(\bar s)=\overline{D(s)}$; in particular $D(1)$ is real.
The function $G$ in the theorem is the real sum of prime-ideal ramps
plus $xD(1)$.  For $\Rea s>2$, absolute convergence and
\[
 \int_n^\infty\left(\frac{x}{n}-1\right)x^{-s-1}\,dx
       =\frac{n^{-s}}{s(s-1)}
\]
give the exact Mellin identity
\begin{equation}\label{eq:mellin}
 \Phi(s):=\int_1^\infty G(x)x^{-s-1}\,dx
       =\frac{sD(1)-D(s)}{s(s-1)}.
\end{equation}
This formula also provides the meromorphic continuation in $\Rea s>0$.

Let $Z$ be the union of the zero sets of the $L(s,\lambda_j)$ and
$L(s,\bar\lambda_j)$ in that half-plane.  The set is discrete because
these finitely many entire functions are not identically zero.  By
Lemma~\ref{lem:one}, every element of $Z$ lies in $0<\Rea s<1$.
The possible pole at one in \eqref{eq:mellin} is removable, since the
numerator vanishes there.  At $\rho\in Z$ put
\[
 M_\rho=\frac12\sum_{j=1}^r a_j
      \bigl(\ord_\rho L(s,\lambda_j)+
                         \ord_\rho L(s,\bar\lambda_j)\bigr)>0.
\]
Then $\Phi$ has a simple pole at $\rho$ with residue
\begin{equation}\label{eq:residue}
             \frac{M_\rho}{\rho(1-\rho)}.
\end{equation}
It is nonzero, and it is positive real if $\rho$ is real.  This also
proves that overlapping zero sets cause no cancellation.

Extend $G(x)$ to $0<x<1$ by $xD(1)$ and put $g(u)=G(e^u)$ for every
$u\in\mathbb R$.  It is continuous: every summand of a ramp vanishes
at its starting point, and the sum is locally finite.  To check its
growth, group prime-ideal powers according to their rational norm.  For
each rational prime power $p^m$, the absolute value of the combined
logarithmic-derivative coefficient is at most $2\log p$ for each
character.  This follows by summing the residue degrees of the prime
ideals above $p$, which is at most two.  Thus
\[
 |G(x)|\leq C_1x\sum_{n\leq x}\frac{\log n}{n}+|D(1)|x
       \leq C_2x(1+\log^2x)\leq C_3x^2\quad(x\geq1).
\]
In particular $|g(u)|\leq C_3e^{2u}$ for $u\geq0$.  With $x=e^u$,
\eqref{eq:mellin} is the Laplace transform
$\Phi(s)=\int_0^\infty g(u)e^{-su}\,du$ for $\Rea s>2$.
We have verified measurability, exponential growth, the Laplace identity,
analyticity off $Z$, isolated poles, and the residue conditions in
Alp\"oge's abstract theorem.

Fix $\theta\in(1/2,1)$ and put $\delta=\delta_\theta$.
The assumed windows become $[U,(1+\delta)U]$, where $U=\log Y$.
On such a window take $R(u)=|C_\theta|e^{\theta u}$.
Then $0\leq R(u)$, $g(u)\leq R(u)$, and, for every $\eta>0$,
\[
 \int_U^{(1+\delta)U}R(u)e^{-\theta u}\,du
       =|C_\theta|\delta U\leq e^{\eta U}
\]
once $U$ is sufficiently large.  Arbitrarily late original windows
therefore verify the complete one-sided window hypothesis of
\cite[\texttt{Deduction.landau\_window}]{AlpogeWindow}.
That theorem gives $\Rea\rho\leq\theta$ for every $\rho\in Z$.
Letting $\theta\downarrow1/2$ excludes all zeros to the right of
$1/2$.  The functional equation, with the conjugate characters already
included, excludes zeros to its left in the critical strip.  Its gamma
factors have neither zeros nor poles in that strip because the angular
types are nonzero.  This proves the theorem.
\end{proof}

\begin{proof}[Proof of Corollary~\ref{cor:three}]
The first assertion is Theorem~\ref{thm:window} with three weights equal
to one and the unitary shifts in \eqref{eq:targets}.
For a prime $\ell$, use the $\ell$-adic realization $\psi_{36,\ell}$ of the CM character, and write $V=\operatorname{Ind}_{G_K}^{G_\Q}\psi_{36,\ell}$.  In the following representation calculation, write $\psi$ for this realization and $\psi^c$ for its Galois conjugate.
On restriction to $G_K$, the four cubic monomials split into the two
pairs with characters
\[
 (\psi^3,(\psi^c)^3),\qquad
 (\psi^2\psi^c,\psi(\psi^c)^2).
\]
The two pairs are stable under $G_\Q$ and exchanged within each pair by
complex conjugation.  Since $\psi\psi^c$ is the cyclotomic
character on $G_K$, this gives
\[
 \operatorname{Sym}^3V\simeq
 \operatorname{Ind}_{G_K}^{G_\Q}(\psi^3)\oplus V(\mathrm{cyc}).
\]
The possible quadratic extension-character twist of the second
summand does not change its isomorphism class, since
$V\otimes\chi_K\simeq V$.  Hence the full finite $L$-function is
\[
 L(\operatorname{Sym}^3V,s)
          =L(s,\psi_{36}^{3})L(s-1,\psi_{36}).
\]
Both factors are entire.  Their critical lines are both $\Rea s=2$
in this product, so their Riemann hypotheses imply that of the product.
This factorization uses the full representations and therefore includes
their ramified local factors.
\end{proof}

\section{Where the extension still requires arithmetic}

This section records which assertions of the proposed extension are
proved and where its unproved inequality enters.

The abstract Landau theorem is already formulated for a general
Laplace transform.  The Dirichlet-specific wrapper takes a finite family
of Dirichlet characters and a centered sum over rational integers
\cite[\texttt{CharFamily}, \texttt{Gchi},
\texttt{zeroFree\_of\_rpow\_window}]{AlpogeCode,AlpogeWindow}.
Theorem~\ref{thm:window} transports this analytic step to the desired
Hecke characters.  The arithmetic passage producing the windows,
\cite[\texttt{windowInequality\_of\_artinInputs} and
\texttt{level\_gm\_le\_on\_family\_scales}]{AlpogeArithmetic}, uses
rational residue-class data and the characters occurring in its level.
For positive infinity type, it must instead control the angular weights
of prime ideals in \eqref{eq:Gdef} with precisely the centering term there.
No such comparison is proved here.

\begin{proposition}[An obstruction to a fixed finite Dirichlet family]
\label{prop:finite-family}
Let $\lambda_1,\lambda_2,\lambda_3$ be the characters in
\eqref{eq:targets}, and let $a_1,a_2,a_3>0$.  The logarithmic derivative
\[
 -\sum_{j=1}^3a_j\frac{L'(s,\lambda_j)}{L(s,\lambda_j)},
\]
with ideal terms collected according to their rational norms, is not a
fixed finite complex linear combination of logarithmic derivatives of
Dirichlet $L$-functions.  The assertion remains true after finitely many
rational Euler factors are deleted on either side.
\end{proposition}

\begin{proof}
Suppose such an identity holds in the half-plane of absolute
convergence.  Choose a split rational prime $p$ outside the deleted
factors, the three conductors, and the conductors of the finitely many
Dirichlet characters.  Choose a prime $\mathfrak p$ above $p$, and put
$u=\psi_{36}(\mathfrak p)/\sqrt p$.  The other prime has value $u^{-1}$
in unitary normalization.  By Proposition~\ref{prop:twist}, for some
$\zeta\in\mu_3$ the three characters at $\mathfrak p$ have values
$u,\zeta u,u^3$.

None of $u,\zeta u,u^3$ is a root of unity.  Indeed, choose a generator $\alpha$ of $\mathfrak p$.  Infinity type one gives $\psi_{36}(\mathfrak p)=\xi\alpha$ for a root of unity $\xi$, and $\alpha\bar\alpha=p$.  If $u^M=1$, choose a positive integer $T$ with $\xi^T=1$.  Then $\alpha^{2MT}=p^{MT}$.  The valuations at $\mathfrak p$ of the two sides are $2MT$ and $MT$, a contradiction.  The same argument applies to
every nonzero power of $u$ and to multiplication by $\zeta$.

Uniqueness of absolutely convergent Dirichlet series identifies the
coefficient of every $p^n$.  Dividing this coefficient by $\log p$ gives
on the Hecke side
\[
 a_1(u^n+u^{-n})+
 a_2((\zeta u)^n+(\zeta^{-1}u^{-1})^n)+
 a_3(u^{3n}+u^{-3n}).
\]
On the proposed Dirichlet side it is a finite linear combination of
$\chi(p)^n$, whose bases $\chi(p)$ are roots of unity.  Collect equal
bases on each side.  On the Hecke side each coefficient is a sum of
positive numbers and therefore nonzero, and each base is not a root
of unity.  The distinct sequences $z^n$, for finitely many nonzero
complex numbers $z$, are linearly independent: evaluate at the first
as many positive integers as there are bases, and use the nonzero
Vandermonde determinant.  This contradicts the coefficient identities.
There are arbitrarily many split primes, so deleting finitely many
Euler factors cannot avoid the argument.
\end{proof}

At $p=7$ the first character already illustrates the obstruction.
Its unitary Frobenius trace is $-4/\sqrt7$; its trace at the second
power is $2/7$, and its trace at the third power is $20/(7\sqrt7)$.
The entire prime-power sequence is not periodic.  In contrast, a
fixed finite combination of Dirichlet character values at powers of
seven is periodic.

The proposition rules out an exact substitution into a fixed finite
character family.  It does not rule out families varying with the
scale, or a new arithmetic comparison incorporating the angular
character.  Such an approximation must also control its error and the
conductor and rank costs throughout the long windows.

There are three elementary restrictions on proposed replacements.
First, a nonzero infinity type cannot factor through the finite unit
colour: Proposition~\ref{prop:colour} gives infinitely many different
types with the same colour.  Second, a trace identity at each finite
field does not bound a sum of those traces over different prime ideals.
Third, the pointwise Gauss--Jacobi identity
$g(\chi)g(\phi)=J(\chi,\phi)g(\chi\phi)$ does not transfer a
cancellation estimate for any one of its Gauss-sum factors to a
cancellation estimate for $J$.  Averages of quotients require their
joint phases.  Even two sequences each summing to zero may have
pointwise quotient identically one.  The Jacobi identity identifies the
desired coefficient; it leaves the requisite average to be bounded.

The order-six residue action on a rank-twenty-four lattice is a genuine
finite representation, but an identification of that representation
with the arithmetic family in the prime-ideal inequality is additional
mathematics.  The centre of the umbral group cannot provide its cubic
part, by Proposition~\ref{prop:colour}.

\subsection{The rank of the theta-lift family}

The mention of a theta-series index in step~18 of Alp\"oge's
announcement \cite[p.~4]{Alpoge} does not define a particular rank-twenty-four
lattice or a root-lattice glue index.  There is also a direct rank
obstruction to identifying the growing chain itself with a fixed
Niemeier lattice.  The formal size conditions require, with
$x>1$ and $\varepsilon>0$,
\[
 k\geq x^{9/2+10\varepsilon-\varepsilon/100},
 \qquad N(j)=k+j(t+1),
\]
and the family has arbitrarily large values of $x$
\cite[\texttt{SizeK}, \texttt{Ranks.N}, \texttt{FamCore.Win\_unb}]{AlpogeRanks}.
The exponent is positive, so $k$ and every $N(j)$ are unbounded.
A lattice whose rank is one of these numbers cannot be a fixed
Niemeier lattice for all large $x$; even its rank over the coefficient
field eventually exceeds twenty-four.  Extracting a special fixed-rank
piece, or identifying a different lattice in step~18, would require an
explicit construction.  Neither a discriminant nor an index follows
from the four-page outline alone.

There is also a normalization issue even for a fixed auxiliary lattice.
Let $L$ be free of rank $n$ over $\mathcal O$, let $h$ be a positive
Hermitian form linear in its first variable with matrix $H$ in an
$\mathcal O$-basis, and put $B=\Tr_{K/\Q}h$.  Define
$L^\#=\{v:h(v,L)\subset\mathcal O\}$ and let
$\mathfrak D=(\sqrt{-3})$ be the different.  Then
\begin{equation}\label{eq:trace-lattice}
 \det_{\Z}(L,B)=3^n(\det H)^2,\qquad
 L^{\vee_B}=\mathfrak D^{-1}L^\#.
\end{equation}
For the standard rank-one Hermitian lattice, the trace Gram matrix
in the basis $1,\omega$ is
$\left(\begin{smallmatrix}2&-1\\-1&2\end{smallmatrix}\right)$,
of determinant three.  A complex change of variables with
$H=C C^*$ has real determinant $|\det C|^2=\det H$, and a
Gram determinant changes by its square.  This proves the first
formula.  For the second, testing $\Tr h(v,a\ell)$ for every
$a\in\mathcal O$ and $\ell\in L$ says precisely that
$h(v,L)\subset\mathfrak D^{-1}$, which is equivalent to
$\sqrt{-3}\,v\in L^\#$.

Thus an underlying Niemeier lattice requires $n=12$, an integral
even trace form, and $L^\#=\mathfrak D L$; its Hermitian determinant
in this convention is $3^{-6}$.  A Hermitian integral lattice with
$\det H=1$ instead has trace determinant $3^n$, so Hermitian
self-duality is not the required condition.  If the trace form is
scaled by a positive rational $c$, the determinant becomes
$c^{2n}3^n(\det H)^2$ and the dual becomes
$(c\mathfrak D)^{-1}L^\#$.  These elementary formulas were
independently derived in the lattice companion calculation.  No
choice of the scale, an integral structure, or a rank-twelve
summand is specified by step~18 itself.

\section{A shifted auxiliary factor in the sextic probe}

There is a recent extension of the cubic-theta method to continuous
unitary Hecke characters.  Gao's Theorem~1 gives the boundary
$(777-\sqrt{1081})/852$, subject to the analytic inputs explicitly listed
there \cite{Gao}.  Liu's Theorem~7.33 states the stronger boundary
\[
 .8733326374459049352<\theta_{\rm G}<.8733326374459051924
\]
for every number field and every continuous unitary Hecke character
\cite{LiuMixed}.  These results include nonzero infinity type; their
intermediate constants may depend on that fixed type.  Neither states
RH.  We have checked the statements, their cited dependencies, and the
local calculation used below, but have not independently audited their
full analytic proofs or the numerical certificates named in Liu's text.
They are not premises of Theorem~\ref{thm:window}.

The following calculation sharpens one particular interface in that
method.  Hu's principal Euler factor \cite[Section~6.2, Lemma~6.1]{Hu}
has the form used in the proposition.  Its original absolutely convergent
product is defined to the right of $2/3$.

\begin{proposition}[Exact continuation of the principal factor]
\label{prop:auxiliary}
Let $F$ be a number field, let $\psi,\chi$ be unitary Hecke characters,
and let $S$ contain their ramified finite places.  For
$q=\N\mathfrak p$ and $\mathfrak p\notin S$, put
\[
 D_{\mathfrak p}=\psi(\mathfrak p)q^{-s},\qquad
 E_{\mathfrak p}=\chi(\mathfrak p)q^{3-6s},\qquad
 H_{\mathfrak p}(s)=
 \frac{(1-q^{-2})\{1-q^{-1}E_{\mathfrak p}
                     -(1-q^{-1})D_{\mathfrak p}\}}
 {(1-E_{\mathfrak p})(1-D_{\mathfrak p})}.
\]
The product $H(s)=\prod_{\mathfrak p\notin S}H_{\mathfrak p}(s)$,
initially for $\Rea s>2/3$, has the meromorphic continuation
\begin{equation}\label{eq:Hfactor}
 H(s)=\zeta_F^S(2)^{-1}L_F^S(6s-3,\chi)J(s)
 \qquad(\Rea s>1/2),
\end{equation}
where $J$ is holomorphic and nowhere zero throughout that half-plane.
Consequently a zero of $L_F^S(s,\psi)$ at $\rho$ with
$\Rea\rho>1/2$ can be canceled in $H(s)/L_F^S(s,\psi)$ only if
$L_F^S(6\rho-3,\chi)=0$, with at least the same multiplicity.
\end{proposition}

\begin{proof}
Take
\[
 J_{\mathfrak p}(s)=
 \frac{1-q^{-1}E_{\mathfrak p}-(1-q^{-1})D_{\mathfrak p}}
      {1-D_{\mathfrak p}}
 =1+\frac{q^{-1}(D_{\mathfrak p}-E_{\mathfrak p})}
            {1-D_{\mathfrak p}}.
\]
On every compact subset of $\Rea s>1/2$, with minimum real part
$\sigma_0>1/2$, the denominator is bounded away from zero and
\[
 |J_{\mathfrak p}(s)-1|
   \ll q^{-1-\sigma_0}+q^{2-6\sigma_0}.
\]
Both prime-ideal sums converge, so the product of these factors
converges normally.  Each factor is nonzero: both $D_{\mathfrak p}$
and $E_{\mathfrak p}$ lie in the open unit disk, and
$q^{-1}E_{\mathfrak p}+(1-q^{-1})D_{\mathfrak p}$ is a strict
convex combination of two such points.  Thus $J$ is holomorphic and
nonzero.  Extracting $(1-E_{\mathfrak p})^{-1}$ and $1-q^{-2}$
proves \eqref{eq:Hfactor} first in the region of absolute convergence
and then by meromorphic continuation.  Since $J(\rho)\ne0$, the
last assertion follows by comparing orders of zeros and poles.  A pole
of the numerator would increase the pole of the quotient and cannot
cancel it.
\end{proof}

\subsection{The additional character for the three CM targets}

Hu's calibration takes a unitary primary-generator character $\Lambda$
and a finite character $G$, with
\[
 \vartheta=\overline\Lambda\,\overline G\psi,
 \qquad\chi=\vartheta^6.
\]
Here $\Lambda$ is the norm pullback of
$\Lambda_K=\psi_{27}\N^{-1/2}$; the quadratic-index part of $G$ has
order dividing two and its cubic-residue part has order dividing three.
Therefore $G^6=1$ and
\begin{equation}\label{eq:auxiliary-character}
 \chi=(\psi/\Lambda)^6.
\end{equation}
These are the arithmetic calibration and the fixed-type orientation
of \cite[Sections~2.2--2.3]{Hu}.  In that construction the period angular
type is the target type minus the type of $\Lambda$, and at least one
of these differences is required to be nonzero.

The working field is $F=\Q(\zeta_{12})=K(i)$ and
$\Lambda=\Lambda_K\circ\N_{F/K}$.  Let a subscript $F$ denote norm
pullback from $K$.  The ratio $\lambda_1/\lambda_2$ has order dividing
three, whereas $\lambda_3=\Lambda_K^3$.
Thus the following admissible orientations give
\[
\begin{array}{c|c|c}
 \text{target over }K&\text{oriented target }\psi\text{ over }F
       &\chi=(\psi/\Lambda)^6\\ \hline
 \lambda_1&\overline{\lambda_{1,F}}&\Lambda^{-12}\\
 \lambda_2&\overline{\lambda_{2,F}}&\Lambda^{-12}\\
 \lambda_3&\lambda_{3,F}&\Lambda^{12}.
\end{array}
\]
For the first two rows the un-conjugated choice would make the period
angular type zero at every complex place.  It would give $\chi=1$,
but does not satisfy the hypothesis of the cited fixed-type reflection.
Using that choice requires a separate treatment of the constant term.

If $\eta$ is the quadratic Hecke character of $F/K$, abelian base
change gives
\[
 L_F(s,\Lambda^{12})
 =L_K(s,\Lambda_K^{12})L_K(s,\Lambda_K^{12}\eta),
\]
and the conjugate identity.  Hence the local continuation introduces
two characters of angular type twelve downstairs.  Neither is a
Dirichlet character.  This identifies additional arithmetic data;
it does not supply the needed bound for the full theta numerator.

\subsection{What cancellation and iteration would require}

Liu's final Mellin step \cite[Section~7.9]{LiuMixed} uses nonvanishing of
$H_\psi$ for $\Rea s>2/3$ and chooses a saving smaller than
$\beta-2/3$.  Proposition~\ref{prop:auxiliary} continues the factor
farther but exposes its actual possible zeros.  It therefore replaces
that absolute-convergence restriction by an explicit cancellation
question; it does not remove the question.

The affine map in this question has fixed point $3/5$, since
$6\beta-3>\beta$ exactly when $\beta>3/5$.  The following conditional
observation makes the scope of such an iteration precise.  Suppose
a family $\mathcal A$ is closed under the map
$T\psi=(\psi/\Lambda)^6$, and let $\beta\le1$ be the supremum of
the real parts of its nontrivial zeros, with $1/2$ included in taking
the supremum.  Suppose, additionally, that for some common $e>0$
every quotient $H_\psi/L_F^S(\cdot,\psi)$ is holomorphic in
$\Rea s>\beta-e$, with the factors of
Proposition~\ref{prop:auxiliary}.  Then $\beta\le3/5$.
Indeed, if $\beta>3/5$, choose a zero $\rho$ with real part
sufficiently close to $\beta$ that
\[
 \Rea\rho>\max\{1/2,\beta-e,(\beta+3)/6\}.
\]
Holomorphy and Proposition~\ref{prop:auxiliary} force a zero of
$L_F(\cdot,T\psi)$ at $6\rho-3$, whose real part is greater than
$\beta$, a contradiction.  Deleted Euler factors have zeros only
on real part zero and do not affect this argument.

Neither the required holomorphy estimate nor closure of the three-target
family has been established.  On angular types the map is
$m\mapsto6(m-1)$ (and conjugation changes the sign).  Starting with
the admissible types above introduces $\pm12$ and then unbounded
absolute types; for $|m|\ge2$ the next absolute type is at least
$6|m|-6\ge3|m|$.  Thus a finite-family argument cannot close this
particular iteration.  Even the conditional conclusion $\beta\le3/5$
does not give RH.  For a hypothetical analogous factor
$L(bs-b/2,T\psi)$ the same calculation has fixed point
$b/(2(b-1))$, which tends to $1/2$ as $b$ increases.  No such family of
probe identities or its uniform global estimates is constructed here.

\section{What remains open}

For the three target characters, the outstanding claim is exactly the
one-sided bound \eqref{eq:window} for the explicitly centered sum
\eqref{eq:Gdef}, or another estimate sufficient to imply it.
The finite-order phase calculation, the cubic-twist identification,
and the analytic implication are established here.  The uniform
arithmetic comparison with the metaplectic theta construction remains
to be made.  Consequently this note proves none of the three requested
Hecke Riemann hypotheses.

\paragraph{Method and records.}
The mathematics and text were produced by GPT-6 Astra in ChatGPT Pro
mode.  Another instance of the same model independently re-proved
Theorem~\ref{thm:window} from its statement, granting the cited abstract
Landau theorem and the classical Hecke continuation and functional
equation.  The conductor and symmetric-cube local calculations were
also independently checked by the same model in the companion work.
These are mathematical checks, not a new Lean certification, and no
human has checked this text.  Our text and original programs are
dedicated to the public domain under
\href{https://creativecommons.org/publicdomain/zero/1.0/}{CC0 1.0}.
The workbench is \url{\workbenchurl}; the permanent address of its
records is \url{\recordurl}.

\begin{thebibliography}{9}
\raggedright
\bibitem{Alpoge}
L.~Alp\"oge, announcement and remarks on the Riemann hypothesis,
20 November 2026, four pages, especially p.~3, step~3 and footnote~3,
and p.~4, steps~8 and~18.
\url{https://alpo.ge/remark.pdf}.

\bibitem{AlpogeCode}
L.~Alp\"oge, Lean formalization of RH and Dirichlet GRH,
mirror at commit \texttt{7cbf7a0c32fca1badf05931879a37c3f8a2b03b4}.
The declarations
\href{https://github.com/CocoSgt/riemann-hypothesis-lean/blob/7cbf7a0c32fca1badf05931879a37c3f8a2b03b4/src/Statements.lean}{\texttt{DirichletLNonvanishingRightOfHalf}}
and
\href{https://github.com/CocoSgt/riemann-hypothesis-lean/blob/7cbf7a0c32fca1badf05931879a37c3f8a2b03b4/src/Development/AutomorphicArithmetic/DedekindLFunctions/CharFamily.lean}{\texttt{CharFamily}, \texttt{Gchi}},
lines 6--8 and 16--41, respectively.

\bibitem{AlpogeWindow}
L.~Alp\"oge, same formalization,
\href{https://github.com/CocoSgt/riemann-hypothesis-lean/blob/7cbf7a0c32fca1badf05931879a37c3f8a2b03b4/src/Development/AutomorphicArithmetic/DedekindLFunctions/LandauZeroFreeCriterion.lean}{\texttt{PoleData}, \texttt{WindowHypU}},
lines 15--30;
\href{https://github.com/CocoSgt/riemann-hypothesis-lean/blob/7cbf7a0c32fca1badf05931879a37c3f8a2b03b4/src/Development/AutomorphicArithmetic/DedekindLFunctions/LandauTheorems/landauWindow.lean}{\texttt{Deduction.landau\_window}},
lines 83--85, with proof through line 170;
\href{https://github.com/CocoSgt/riemann-hypothesis-lean/blob/7cbf7a0c32fca1badf05931879a37c3f8a2b03b4/src/Development/AutomorphicArithmetic/DedekindLFunctions/LandauTheorems/ZeroFreeOfRpowWindow.lean}{\texttt{zeroFree\_of\_rpow\_window}},
lines 11--14.

\bibitem{AlpogeArithmetic}
L.~Alp\"oge, same formalization,
\href{https://github.com/CocoSgt/riemann-hypothesis-lean/blob/7cbf7a0c32fca1badf05931879a37c3f8a2b03b4/src/Development/AutomorphicArithmetic/PrimitiveCosets/WindowInequalityOfArtinInputs.lean}{\texttt{windowInequality\_of\_artinInputs}},
lines 17--27; and
\href{https://github.com/CocoSgt/riemann-hypothesis-lean/blob/7cbf7a0c32fca1badf05931879a37c3f8a2b03b4/src/Development/AutomorphicArithmetic/PrimitiveCosets/LevelGmLeOnFamilyScales.lean}{\texttt{level\_gm\_le\_on\_family\_scales}},
lines 28--30.

\bibitem{AlpogeRanks}
L.~Alp\"oge, same formalization,
\href{https://github.com/CocoSgt/riemann-hypothesis-lean/blob/7cbf7a0c32fca1badf05931879a37c3f8a2b03b4/src/Formalism/SharpRestriction/ElevenLe.lean}{\texttt{SizeK}}, lines 31--33;
\href{https://github.com/CocoSgt/riemann-hypothesis-lean/blob/7cbf7a0c32fca1badf05931879a37c3f8a2b03b4/src/Formalism/LocalRingCardinality/Params.lean}{\texttt{Ranks.N}}, line 36;
\href{https://github.com/CocoSgt/riemann-hypothesis-lean/blob/7cbf7a0c32fca1badf05931879a37c3f8a2b03b4/src/Formalism/SharpRestriction/FamCore.lean}{\texttt{FamCore.Win\_unb}}, line 79.

\bibitem{Prasad}
D.~Prasad, \emph{Lectures on Tate's thesis}, ICTP lecture notes,
Theorem~2, p.~3; Corollary~2, p.~9; Theorem~4 and proof, pp.~10--11.
\url{https://www.math.iitb.ac.in/~dprasad/ictptate.pdf}.

\bibitem{Hu}
B.~Hu, \emph{A seven-eighths zero-free half-plane for unitary Weil
$L$-functions}, Hexagon 2610.00216v2, 2026.
\url{https://hexagonmath.org/2610.00216v2}.
Mathematical locators use the pinned
\href{https://github.com/Rogerhu12/hecke-seven-eighths/blob/edbecbec14e1a2efd50d1f0606b475fa5eab88f8/weil_seven_eighths_revised.tex}{source revision \texttt{edbecbec}}:
arithmetic calibration and infinity-type orientation,
lines 442--635; principal Euler factors and Lemma~6.1,
lines 3444--3515.

\bibitem{OpenAIQRH}
OpenAI, \emph{The Quasi-Riemann Hypothesis: A Zero-Free Half-Plane
$\Rea s>7/8$}, 30 September 2026, 199-page edition,
Theorem~1.1, p.~4, and proof overview, pp.~5--7.
\href{https://github.com/openai/math/blob/main/preprints/The-Quasi-Riemann-Hypothesis-September-30-2026/paper.pdf}{Dated primary release and paper}.

\bibitem{Gao}
Z.~Gao, \emph{Unitary character moments and a zero-free half-plane
for unitary Weil $L$-functions}, Hexagon 2611.00334v1,
11 November 2026, Theorems~1--2 and the analytic input list after
Theorem~1; source lines 77--109 and 175--204.
\url{https://hexagonmath.org/2611.00334v1}.

\bibitem{LiuMixed}
B.~Liu, \emph{Mixed character moments and a zero-free half-plane
for unitary Hecke $L$-functions}, Hexagon 2611.00741v1,
20 November 2026, Section~7.9 and Theorem~7.33 in Section~7.18;
source labels \texttt{subsec:full-physical-return} and
\texttt{thm:enhanced-field}.
\url{https://hexagonmath.org/2611.00741v1}.

\bibitem{LeanCommunity}
Lean FRO and Mathlib Initiative, \emph{Lean verification of the Riemann
Hypothesis}, joint statement, 21 November 2026, verification criteria
1--4.
\url{https://leanprover-community.github.io/blog/posts/riemann-hypothesis/}.

\bibitem{MITVerification}
A.~Sutherland, \emph{Verification files for the Lean proof of the Riemann
Hypothesis}, summary and five-target verification table, updated
24 November 2026, 05:25 UTC.
\url{https://math.mit.edu/~drew/grh/}.
\end{thebibliography}
\end{document}
